Autonomous visual navigation for small drones sits at the intersection
of several conflicting constraints: low mass, low power, real-time
perception, onboard decision making, and enough software flexibility to
support rapid iteration. In practice, many platforms solve only part of
this problem. Very small aerial vehicles are easy to fly indoors and
safe to test in tight environments, but they often lack the compute and
memory budget required by modern object detection pipelines. At the
opposite end, larger embedded AI platforms such as Coral Dev Board or
NVIDIA Jetson-class systems provide significantly more resources, but
they also bring higher current consumption, larger batteries, more
weight, and a form factor that is difficult to justify for lightweight
micro-UAV missions.

This work starts from an open-source platform with clear latent
potential, namely Coral Micro, and treats it not as a finished product
but as a promising technical base that can be pushed toward a much
higher level of systems maturity. The research direction is therefore
not limited to adding isolated features. Instead, it aims to move an
accessible open hardware and open software project toward an
industrialized embedded-AI stack with production-oriented performance
characteristics, stronger runtime guarantees, richer sensing, and a more
coherent control and telemetry model.

The starting point of this work was the need for a platform that remains
physically lightweight and energy-efficient, while still exploiting as
much as possible the inference capability of the EdgeTPU and the
real-time performance of the MCU. The original minimalist EdgeTPU-based
board configuration was attractive from the perspective of size and
power, but its standard low-resolution camera workflow was not
sufficient for the class of missions targeted here. In particular, a
320x240 sensing regime is too restrictive when the goal is to support
richer object detection, stricter scene understanding, and
higher-confidence flight decisions in inhabited or cluttered spaces.
This became even more limiting when considering modern YOLO-style
detectors, which benefit from larger spatial inputs and from a sensor
pipeline capable of preserving more detail before resizing and
quantization.

The objective of SentAI is therefore to define a new point in the design
space: a board-level and software-level architecture that preserves the
low-weight, low-power philosophy of minimalist embedded AI hardware, but
extends it toward practical autonomous drone operation. The platform
couples an MCU and an EdgeTPU with two 5 Mpx cameras and an onboard IMU,
and exposes the whole sensing, inference, communication, and control
stack through a MicroPython environment. This makes it possible to write
high-level mission logic in Python while keeping timing-critical
operations in the underlying real-time operating system and accelerator
path. The result is not only a development framework, but also a systems
study of how far one can push a compact embedded AI board toward real
airborne autonomy and toward a level of robustness closer to
production-grade embedded deployment.

The work explicitly targets two flight platforms. The first is Bitcraze
Crazyflie, used as an indoor experimental platform because it is agile,
lightweight, easy to maneuver, and well suited for repeated
safety-conscious testing in the laboratory. The second is a PX4-based
autopilot stack accessed through MAVLink, representing the path toward a
more professional and transferable deployment model. The intention is to
provide a common high-level autonomy layer across both targets, so that
mission scripts, perception logic, and communication flows can be
developed once and adapted across indoor research flights and more
operational MAVLink-centered scenarios.

From a vision-systems perspective, the work is driven by the observation
that model accuracy alone is not enough. Once YOLO-class detectors are
pushed toward larger inputs, for example 640-class or even workflows
derived from native 1280-wide camera imagery, the dominant issues become
memory layout, transfer scheduling, tensor preparation cost, and buffer
ownership between the camera subsystem, the MCU, and the EdgeTPU. For
this reason, the contribution is as much about systems engineering as it
is about AI. The implementation analyzes the full visual path from image
acquisition at the sensor, through DMA-backed capture into RAM, through
hardware pixel conversion and scaling, and finally into the EdgeTPU
input memory. These stages are pipelined in parallel under the RTOS so
that camera capture and image preparation can overlap with inference,
thereby maximizing frame rate while minimizing idle time and unnecessary
copies.

An important architectural extension is the addition of a dual-camera
subsystem. Instead of treating perception as a single fixed viewpoint
problem, SentAI supports two 5 Mpx sensors that may be mounted in
parallel or perpendicular configurations. This enables distinct
operational modes: forward-looking navigation together with nadir
inspection, redundant views for indoor localization, or viewpoint
switching according to the phase of the mission. The camera-switching
mechanism is designed as a first-class runtime concept, not as an
afterthought, and is supported by explicit per-camera geometry, frame
sequencing, and clean post-switch acquisition logic. Together with
IMU-aware tracking and ground-plane projection, this creates the basis
for drone-centric perception that is aware of mounting geometry and
flight context rather than only raw image content.

The work also extends beyond the runtime itself. In addition to dynamic
model loading and multi-model switching, the project assumes the
availability of a custom YOLO detector trained on an open synthetic
dataset built for low-altitude drone navigation in inhabited spaces. The
dataset contains views captured at nadir and at approximately 45
degrees, and is intended as a reusable pre-training or adaptation
resource for future systems that must obey stricter navigation and
avoidance logic in populated environments, including scenarios informed
by European drone categories such as A0-A3 and A2. This
dataset-and-model layer complements the embedded runtime by providing a
path from open training assets to open onboard deployment.

In this sense, the present work should be read as an industrialization
effort. Many of the underlying techniques are known in isolation:
embedded Python, edge inference, object detection, MAVLink
communication, or visual tracking. The novelty lies in bringing them
together on a highly constrained airborne computing platform and
resolving the practical bottlenecks that appear only during
implementation. These include memory fragmentation, RAM placement
between internal memory and external SDRAM, large tensor staging,
synchronization between RTOS tasks, camera switching latency, and the
challenge of sustaining high throughput without losing the low-power
character of the platform. The broader contribution is thus to
demonstrate how an open-source prototype platform can be evolved into a
system with much stronger production-oriented behavior, where
throughput, determinism, integration discipline, and operational
usability become design targets in their own right.

The main contributions of the work can be summarized as follows:

\begin{enumerate}
\def\labelenumi{\arabic{enumi}.}
\tightlist
\item
  An open-source MicroPython runtime for a lightweight EdgeTPU-enabled
  MCU platform intended for high-level autonomous drone control.
\item
  A dual-target flight-control concept spanning Bitcraze Crazyflie for
  indoor experimentation and PX4/MAVLink for professional autopilot
  integration.
\item
  A dual-camera 5 Mpx sensing architecture with onboard IMU support,
  runtime camera switching, per-camera geometry, and viewpoint-aware
  operation.
\item
  A real-time perception pipeline that studies and optimizes the
  end-to-end path from image sensor to MCU RAM to EdgeTPU memory, using
  parallel RTOS execution and double-buffered staging to maximize frame
  rate.
\item
  A memory-aware deployment strategy for YOLO-class detectors and other
  larger vision models on an MCU-constrained platform.
\item
  A multi-model workflow with dynamic model loading, context-dependent
  switching, tracking, projection to the ground plane, and telemetry
  integration.
\item
  An open synthetic aerial dataset and a custom YOLO model intended as a
  starting point for regulation-aware drone navigation in inhabited
  environments.
\end{enumerate}

The rest of the paper is organized as follows. The next section reviews
related work in autonomous nano-UAV navigation, embedded deep learning,
YOLO-based detection, MicroPython-enabled embedded control, and
accelerator-centric edge AI. The following sections present the hardware
and software architecture, the real-time memory and processing pipeline,
the communication and flight-control abstractions, and the experimental
perspective enabled by the platform.
